\begin{table}[htbp]
\centering
\caption{稳健性检验：OLS与DML估计对比}
\label{tab:ols_vs_dml}
\begin{threeparttable}
\begin{tabular}{lcccc}
\toprule
 & \multicolumn{2}{c}{传统OLS} & \multicolumn{2}{c}{DML-CRE} \\
\cmidrule(lr){2-3}\cmidrule(lr){4-5}
 & (1) OLS-15X & (2) OLS-28X & (3) DML-15X & (4) DML-28X \\
\midrule
DU\_kw & $-0.0034$ & $-0.0026$ & $-0.0025^{***}$ & $-0.0033^{***}$ \\
 & $(0.0010)$ & $(0.0010)$ & $(0.0003)$ & $(0.0003)$ \\
\midrule
$N$ & 40,135 & 40,135 & 45,506 & 40,135 \\
控制变量 & 15个 & 28个 & 15个+CRE & 28个+CRE \\
企业FE & \checkmark & \checkmark & Mundlak & Mundlak \\
年份FE & \checkmark & \checkmark & 虚拟变量 & 虚拟变量 \\
\bottomrule
\end{tabular}
\begin{tablenotes}
\small
\item 注：OLS列使用双向固定效应 (企业+年份)，行业×年份聚类标准误。
DML列使用CRE-Mundlak装置，White HC稳健标准误。
***、**、* 分别表示在1\%、5\%、10\%水平显著。
\end{tablenotes}
\end{threeparttable}
\end{table}
